%% ═══════════════════════════════════════════════════════════════════════════
%% Study Notes Template — Polished Edition
%% Fixes applied:
%%   · tcolorbox drop shadow corrected
%%   · AutoFakeBold/AutoFakeSlant replaced with explicit font declarations
%%   · \inlinecode uses \colorbox instead of TikZ
%%   · Redundant color aliases removed
%%   · Widow/orphan protection added
%%   · \parskip set with explicit length
%%   · lstset gets a default language
%%   · \titleformat section guard removed; use \section* when unnumbered
%%   · Subject/title parameterised via commands
%%   · TOC styling added
%% ═══════════════════════════════════════════════════════════════════════════

\documentclass[12pt,a4paper]{article}

%% ─── Packages ────────────────────────────────────────────────────────────────
\usepackage{fontspec}

%% If your font has real bold/italic cuts, declare them explicitly:
\setmainfont{Anthropic-serif}[Path=./, Extension=.ttf]
%% \setmainfont{Anthropic-serif}[
%%     Path        = ./,
%%     Extension   = .ttf,
%%     BoldFont    = Anthropic-serif-Bold,       %% replace with actual filename
%%     ItalicFont  = Anthropic-serif-Italic,     %% replace with actual filename
%%     BoldItalicFont = Anthropic-serif-BoldItalic
%% ]
%% Fallback if only one weight exists — remove the lines above and uncomment:
%% \setmainfont{Anthropic-serif}[Path=./, Extension=.ttf]

\setsansfont{Poppins}
\setmonofont{DejaVu Sans Mono}[Scale=MatchLowercase]  %% or any monospace you have

\usepackage[margin=2.2cm, top=2.5cm, bottom=2.5cm]{geometry}
\usepackage[table]{xcolor}
\usepackage{tikz}
\usetikzlibrary{shapes.geometric, arrows.meta, positioning, fit, backgrounds,
                decorations.pathreplacing, calc}
\usepackage{listings}
\usepackage{tcolorbox}
\tcbuselibrary{listings, skins, breakable, raster}
\usepackage{fancyhdr}
\usepackage{titlesec}
\usepackage{titletoc}
\usepackage{enumitem}
\usepackage{booktabs}
\usepackage{tabularx}
\usepackage{array}
\usepackage{microtype}
\usepackage{hyperref}
\usepackage{amssymb}
\usepackage{graphicx}
\usepackage{amsmath}
\usepackage{parskip}
\usepackage{colortbl}

%% ─── Document metadata (edit here only) ─────────────────────────────────────
\newcommand{\docsubject}{Computer Engineering Theory}
\newcommand{\docshortsubject}{CET}
\newcommand{\doctitle}{Study Notes}

%% ─── Widow / orphan protection ───────────────────────────────────────────────
\widowpenalty=10000
\clubpenalty=10000

%% ─── Paragraph spacing ───────────────────────────────────────────────────────
\setlength{\parskip}{0.5em plus 0.2em minus 0.1em}
\setlength{\parindent}{0pt}

%% ─── Table geometry ──────────────────────────────────────────────────────────
\setlength{\aboverulesep}{0pt}
\setlength{\belowrulesep}{0pt}
\setlength{\extrarowheight}{0.75ex}
\renewcommand{\arraystretch}{1.5}
\setlength{\arrayrulewidth}{0.3pt}

\newcolumntype{L}[1]{>{\raggedright\arraybackslash}p{#1}}
\newcolumntype{C}[1]{>{\centering\arraybackslash}p{#1}}
\newcolumntype{R}[1]{>{\raggedleft\arraybackslash}p{#1}}

%% ─── Brand Colors (Anthropic palette — light mode) ───────────────────────────
\definecolor{AnthDark}{HTML}{141413}
\definecolor{AnthLight}{HTML}{faf9f5}
\definecolor{AnthMidGray}{HTML}{b0aea5}
\definecolor{AnthLightGray}{HTML}{e8e6dc}
\definecolor{AnthOrange}{HTML}{d97757}
\definecolor{AnthBlue}{HTML}{6a9bcc}
\definecolor{AnthGreen}{HTML}{788c5d}
\definecolor{AnthCodeBG}{HTML}{e8e6dc}

\arrayrulecolor{AnthMidGray!50}

%% ─── Hyperref ────────────────────────────────────────────────────────────────
\hypersetup{
  colorlinks  = true,
  linkcolor   = AnthBlue,
  urlcolor    = AnthOrange,
  citecolor   = AnthGreen,
  pdftitle    = {\doctitle{} · \docshortsubject},
  pdfauthor   = {Sudheesh},
}

%% ─── Page style ──────────────────────────────────────────────────────────────
\pagestyle{fancy}
\fancyhf{}
\fancyhead[L]{\small\color{AnthMidGray}\textbf{\doctitle} · \docshortsubject}
\fancyhead[R]{\small\color{AnthMidGray}\thepage}
\renewcommand{\headrulewidth}{0.4pt}
\renewcommand{\headrule}{%
  \hbox to\headwidth{\color{AnthOrange}\leaders\hrule height \headrulewidth\hfill}}
\renewcommand{\footrulewidth}{0pt}

%% ─── Section formatting ──────────────────────────────────────────────────────
\titleformat{\section}
  {\Large\bfseries\sffamily\color{AnthDark}}
  {\colorbox{AnthOrange}{\makebox[0.5cm][c]{\color{white}\thesection}}\quad}
  {0em}
  {}
  [\vspace{-0.5ex}\color{AnthLightGray}\hrule\vspace{0.5ex}]

\titleformat{\subsection}
  {\large\bfseries\sffamily\color{AnthDark}}
  {\color{AnthBlue}\thesubsection\enspace}
  {0em}{}

\titleformat{\subsubsection}
  {\normalsize\bfseries\sffamily\color{AnthDark}}
  {\color{AnthGreen}\thesubsubsection\enspace}
  {0em}{}

\titlespacing{\section}     {0pt}{1.5em}{0.8em}
\titlespacing{\subsection}  {0pt}{1em}{0.4em}
\titlespacing{\subsubsection}{0pt}{0.8em}{0.3em}

%% ─── TOC styling ─────────────────────────────────────────────────────────────
\contentsmargin{0pt}

\titlecontents{section}
  [0pt]
  {\vspace{4pt}\bfseries\sffamily\color{AnthGreen}}
  {\colorbox{AnthOrange}{\makebox[0.45cm][c]{\color{white}\small\thecontentslabel}}\enspace}
  {}
  {\hfill\color{AnthGreen}\contentspage}

\titlecontents{subsection}
  [1.4em]
  {\small\color{AnthBlue}}
  {\thecontentslabel\enspace}
  {}
  {\titlerule*[6pt]{\tiny\color{AnthMidGray}.}\color{AnthBlue}\contentspage}

\titlecontents{subsubsection}
  [3.2em]
  {\footnotesize\color{AnthBlue}}
  {\thecontentslabel\enspace}
  {}
  {\titlerule*[6pt]{\tiny\color{AnthMidGray}.}\color{AnthBlue}\contentspage}

%% ─── Code Listings ───────────────────────────────────────────────────────────
\lstset{
  language        = Python,           %% sensible default; override per listing
  basicstyle      = \ttfamily\small,
  keywordstyle    = \color{AnthOrange}\bfseries,
  stringstyle     = \color{AnthGreen},
  commentstyle    = \color{AnthMidGray}\itshape,
  numberstyle     = \tiny\color{AnthMidGray},
  numbers         = left,
  stepnumber      = 1,
  numbersep       = 10pt,
  breaklines      = true,
  tabsize         = 4,
  showstringspaces= false,
  captionpos      = b,
  frame           = none,
  xleftmargin     = 18pt,
  xrightmargin    = 6pt,
  aboveskip       = 4pt,
  belowskip       = 4pt,
  keepspaces      = true,
}

%% ─── tcolorbox styles ────────────────────────────────────────────────────────
\tcbset{
  %% Dark code box
  codebox/.style={
    enhanced, breakable,
    colback     = AnthDark,
    colframe    = AnthDark,
    arc=6pt, outer arc=6pt,
    sharp corners=west,
    left=4pt, right=4pt, top=4pt, bottom=4pt,
    boxrule     = 0.6pt,
    colupper    = AnthLight,
    coltitle    = white,
    colbacktitle= AnthDark!80,
    fonttitle   = \small\bfseries\sffamily,
    attach boxed title to top left={yshift=-2mm, xshift=8mm},
    boxed title style={arc=3pt, outer arc=3pt, colback=AnthDark!80, boxrule=0pt, sharp corners=west},
    drop fuzzy shadow={AnthMidGray!40},
    top=8pt,
  },
  %% Neutral note
  notebox/.style={
    enhanced,
    colback     = AnthLightGray!40,
    colframe    = AnthMidGray,
    arc=6pt, outer arc=6pt,
    sharp corners=west,
    left=10pt, right=10pt, top=8pt, bottom=8pt,
    boxrule     = 0.6pt,
    colupper    = AnthDark,
    fonttitle   = \bfseries\sffamily,
    coltitle    = AnthDark,
    colbacktitle= AnthLightGray,
    borderline west={3pt}{0pt}{AnthMidGray},
    drop fuzzy shadow={AnthMidGray!25},
  },
  %% Tip / positive
  tipbox/.style={
    enhanced,
    colback     = AnthGreen!8,
    colframe    = AnthGreen!70,
    arc=6pt, outer arc=6pt,
    sharp corners=west,
    left=10pt, right=10pt, top=8pt, bottom=8pt,
    boxrule     = 0.6pt,
    colupper    = AnthDark,
    fonttitle   = \bfseries\sffamily,
    coltitle    = white,
    colbacktitle= AnthGreen,
    borderline west={3pt}{0pt}{AnthGreen},
    drop fuzzy shadow={AnthGreen!20},
  },
  %% Warning
  warnbox/.style={
    enhanced,
    colback     = AnthOrange!6,
    colframe    = AnthOrange!70,
    arc=6pt, outer arc=6pt,
    sharp corners=west,
    left=10pt, right=10pt, top=8pt, bottom=8pt,
    boxrule     = 0.6pt,
    colupper    = AnthDark,
    fonttitle   = \bfseries\sffamily,
    coltitle    = white,
    colbacktitle= AnthOrange,
    borderline west={3pt}{0pt}{AnthOrange},
    drop fuzzy shadow={AnthOrange!20},
  },
  %% Definition
  defbox/.style={
    enhanced,
    colback     = AnthLight,
    colframe    = AnthLightGray,
    arc=6pt, outer arc=6pt,
    sharp corners=west,
    left=10pt, right=10pt, top=8pt, bottom=8pt,
    boxrule     = 0.6pt,
    colupper    = AnthDark,
    borderline west={3.5pt}{0pt}{AnthOrange},
    drop fuzzy shadow={AnthMidGray!25},
  },
}

%% ─── Inline helpers ──────────────────────────────────────────────────────────

%% Lightweight inline code — no TikZ overhead
\newcommand{\ic}[1]{%
  \mbox{\colorbox{AnthCodeBG}{\footnotesize\ttfamily\textcolor{AnthBlue}{#1}}}%
}

%% Concept label (orange left-rule + bold heading)
\newcommand{\concept}[1]{%
  \noindent\textcolor{AnthGreen}{\rule{3pt}{10pt}}\enspace
  \textbf{\color{AnthDark}#1}\par\vspace{2pt}%
}

%% ─── List style ──────────────────────────────────────────────────────────────
\setlist[itemize,1]{itemsep=2pt, topsep=4pt, parsep=0pt, label={\color{AnthOrange}\small$\blacktriangleright$}}
\setlist[itemize,2]{itemsep=2pt, topsep=4pt, parsep=0pt, label={\color{AnthBlue}\small$\blacktriangleright$}}
\setlist[itemize,3]{itemsep=2pt, topsep=4pt, parsep=0pt, label={\color{AnthGreen}\small$\blacktriangleright$}}

\setlist[enumerate,1]{itemsep=2pt, topsep=4pt, parsep=0pt, label={\color{AnthOrange}\arabic*.}}
\setlist[enumerate,2]{itemsep=2pt, topsep=4pt, parsep=0pt, label={\color{AnthBlue}\alph*)}}
\setlist[enumerate,3]{itemsep=2pt, topsep=4pt, parsep=0pt, label={\color{AnthGreen}\roman*.}}

%% ─── Custom environments (convenience wrappers) ──────────────────────────────
\newenvironment{Note}[1][Note]{%
  \begin{tcolorbox}[notebox, title=#1]}{\end{tcolorbox}}

\newenvironment{Tip}[1][Tip]{%
  \begin{tcolorbox}[tipbox, title=#1]}{\end{tcolorbox}}

\newenvironment{Warn}[1][Warning]{%
  \begin{tcolorbox}[warnbox, title=#1]}{\end{tcolorbox}}

\newenvironment{Def}[1][]{%
  \begin{tcolorbox}[defbox]\ifx&#1&\else\textbf{#1.}\enspace\fi}{\end{tcolorbox}}

%% ─── Title page ──────────────────────────────────────────────────────────────
\newcommand{\maketitlepage}{%
  \begin{titlepage}
    \pagecolor{AnthDark}
    \color{AnthLight}
    \vspace*{4cm}
    \begin{center}
      {\Huge\bfseries\sffamily \doctitle}\\[0.6em]
      {\large\sffamily\color{AnthOrange} \docsubject}\\[2em]
      \color{AnthMidGray}\rule{0.4\textwidth}{0.4pt}\\[1.5em]
      {\small\sffamily\color{AnthMidGray} Compiled \today}
    \end{center}
    \vfill
  \end{titlepage}
  \pagecolor{AnthLight}
  \color{AnthDark}
}

%% ═══════════════════════════════════════════════════════════════════════════
%% Document
%% ═══════════════════════════════════════════════════════════════════════════
\begin{document}

\maketitlepage

%% ─── Table of Contents ───────────────────────────────────────────────────────
\thispagestyle{empty}
{
  \hypersetup{hidelinks}
  \tableofcontents
}
\newpage

%% ─────────────────────────────────────────────────────────────────────────────
\section{Digital Logic Fundamentals}
%% ─────────────────────────────────────────────────────────────────────────────

\subsection{Boolean Algebra}

Boolean algebra operates over the set $\{0, 1\}$ under three primary operations:
\textbf{AND} ($\cdot$), \textbf{OR} ($+$), and \textbf{NOT} ($\overline{\phantom{x}}$).
Every digital circuit can be expressed as a combination of these primitives.

\begin{Def}[De Morgan's Theorems]
  \begin{align}
    \overline{A \cdot B} &= \overline{A} + \overline{B} \\
    \overline{A + B}     &= \overline{A} \cdot \overline{B}
  \end{align}
  These identities let you push a negation \emph{through} a gate by flipping
  AND $\leftrightarrow$ OR — invaluable for converting between NAND/NOR implementations.
\end{Def}

\begin{Note}
  De Morgan's theorems extend naturally to $n$ variables:
  $\overline{A_1 \cdot A_2 \cdots A_n} = \overline{A_1} + \overline{A_2} + \cdots + \overline{A_n}$.
\end{Note}

\subsection{Logic Gates and Truth Tables}

The table below shows the four most common two-input gates. NAND and NOR are
\emph{functionally complete} individually — any Boolean function can be realised
using only NANDs, or only NORs.

\begin{center}
\begin{tabular}{C{0.8cm} C{0.8cm} C{1.1cm} C{1.1cm} C{1.1cm} C{1.1cm}}
  \rowcolor{AnthDark}
  \color{white}\textbf{A} &
  \color{white}\textbf{B} &
  \color{white}\textbf{AND} &
  \color{white}\textbf{OR} &
  \color{white}\textbf{NAND} &
  \color{white}\textbf{NOR} \\
  \rowcolor{AnthLightGray!60} 0 & 0 & 0 & 0 & 1 & 1 \\
  0 & 1 & 0 & 1 & 1 & 0 \\
  \rowcolor{AnthLightGray!60} 1 & 0 & 0 & 1 & 1 & 0 \\
  1 & 1 & 1 & 1 & 0 & 0 \\
\end{tabular}
\end{center}

\begin{Tip}[Implementation tip]
  Prefer NAND-only designs in CMOS: NAND gates are slightly faster and smaller than
  NOR gates in that technology because PMOS transistors (used in pull-up networks) are
  inherently slower than NMOS.
\end{Tip}

\subsection{Karnaugh Maps}

A Karnaugh map (K-map) is a visual method for minimising Boolean expressions.
Groups of $2^n$ adjacent \ic{1}s correspond directly to product terms in the
simplified sum-of-products (SOP) expression.

\concept{Grouping rules}

\begin{enumerate}
  \item Groups must be rectangular and contain a power-of-two number of cells ($1, 2, 4, 8\ldots$).
  \item Every \ic{1} must be covered by at least one group.
  \item Groups may wrap around edges (the map is toroidal).
  \item Prefer the \emph{fewest, largest} groups to minimise literals.
\end{enumerate}

\begin{Warn}[Don't-care conditions]
  Don't-care entries (\ic{X}) may be included in a group if doing so makes it
  larger — but you are \textbf{not obligated} to cover them. Never include an
  \ic{X} in a group if it does not increase the group's size.
\end{Warn}

%% ─────────────────────────────────────────────────────────────────────────────
\section{Computer Architecture}
%% ─────────────────────────────────────────────────────────────────────────────

\subsection{Instruction Set Architectures}

\begin{Def}[ISA]
  An \textbf{Instruction Set Architecture} is the abstract interface between hardware
  and software: it specifies the programmer-visible registers, data types, addressing
  modes, and the encoding of every instruction the processor can execute.
\end{Def}

The two dominant design philosophies are contrasted below.

\begin{center}
\begin{tabularx}{\linewidth}{L{3.2cm} X X}
  \rowcolor{AnthDark}
  \color{white}\textbf{Property} &
  \color{white}\textbf{RISC} &
  \color{white}\textbf{CISC} \\
  \rowcolor{AnthLightGray!60}
  Instruction count & Many simple instructions & Fewer, complex instructions \\
  Memory access     & Load/store only          & Any instruction can access memory \\
  \rowcolor{AnthLightGray!60}
  Register file     & Large (32+)              & Small (8 in x86) \\
  Microcode         & None (typically)         & Extensive \\
  \rowcolor{AnthLightGray!60}
  Examples          & ARM, RISC-V, MIPS        & x86, x86-64 \\
\end{tabularx}
\end{center}

\subsection{Pipelining}

Pipelining overlaps the execution of multiple instructions by dividing the datapath
into discrete stages. A classic five-stage RISC pipeline:

\begin{enumerate}
  \item \textbf{IF} — Instruction Fetch: read instruction from memory at PC.
  \item \textbf{ID} — Instruction Decode / Register Read.
  \item \textbf{EX} — Execute: ALU operation or address calculation.
  \item \textbf{MEM} — Memory access (loads and stores only).
  \item \textbf{WB} — Write Back: result to register file.
\end{enumerate}

\concept{Pipeline hazards}

\begin{itemize}
  \item \textbf{Structural} — two instructions require the same hardware resource simultaneously.
  \item \textbf{Data} — an instruction depends on a result not yet written back (RAW, WAR, WAW).
  \item \textbf{Control} — branch outcome is unknown; fetched instructions may be wrong.
\end{itemize}

\begin{Note}[Forwarding / bypassing]
  Most RAW data hazards are resolved by \emph{forwarding}: routing the output of the
  EX or MEM stage directly back to the EX stage input, without waiting for WB.
  This eliminates the majority of pipeline stalls without compiler intervention.
\end{Note}

\subsubsection{Branch prediction}

Static prediction (always-taken / always-not-taken) achieves roughly 60\% accuracy.
Dynamic predictors maintain a history table; a 2-bit saturating counter scheme
typically reaches 85–90\% accuracy on real workloads.

\begin{Tip}[Loop unrolling]
  Loop unrolling reduces branch overhead \emph{and} exposes more instruction-level
  parallelism to the scheduler — useful when the loop body is small and the iteration
  count is known at compile time.
\end{Tip}

%% ─────────────────────────────────────────────────────────────────────────────
\section{Operating Systems}
%% ─────────────────────────────────────────────────────────────────────────────

\subsection{Process Scheduling}

\begin{Def}[Scheduling criteria]
  Key metrics: \textbf{CPU utilisation}, \textbf{throughput} (processes/unit time),
  \textbf{turnaround time} (submit $\to$ complete), \textbf{waiting time} (time in ready
  queue), and \textbf{response time} (submit $\to$ first response).
\end{Def}

\subsubsection{Completely Fair Scheduler (CFS)}

Linux's default scheduler assigns a \emph{virtual runtime} (\ic{vruntime}) to each
task. The task with the smallest \ic{vruntime} is always scheduled next; it accumulates
\ic{vruntime} at a rate inversely proportional to its \emph{nice} weight. Tasks are
stored in a red-black tree keyed on \ic{vruntime}, giving $O(\log n)$ scheduling
decisions.

\begin{tcolorbox}[codebox, title=Inspecting CFS scheduling weights (Linux)]
\begin{lstlisting}[language=bash]
# View the scheduling policy and priority of a process
chrt -p <PID>

# Inspect vruntime and load weight via /proc
cat /proc/<PID>/sched

# Change the nice value of a running process
renice -n 5 -p <PID>
\end{lstlisting}
\end{tcolorbox}

\subsection{Memory Management}

\subsubsection{Virtual memory and paging}

Each process sees a private, contiguous virtual address space. The MMU translates
virtual addresses to physical addresses via a \emph{page table}; modern systems use
multi-level tables (e.g.\ four levels on x86-64) to avoid storing entries for the
entire virtual space.

\begin{Def}[Translation Lookaside Buffer (TLB)]
  A small, fully-associative cache inside the MMU storing recent virtual-to-physical
  mappings. A TLB \emph{hit} resolves a translation in 1–2 cycles; a \emph{miss}
  requires a page-table walk costing 10–100+ cycles.
\end{Def}

\begin{Warn}[TLB shootdowns in SMP]
  When a mapping is modified on one core, all other cores must invalidate their
  TLB entries for that page. This inter-processor interrupt (IPI) broadcast —
  the \emph{TLB shootdown} — is a significant scalability bottleneck on
  many-core systems under high \ic{munmap}/\ic{mmap} rates.
\end{Warn}

\subsubsection{Page replacement policies}

\begin{center}
\begin{tabularx}{\linewidth}{L{2.4cm} X L{2.8cm}}
  \rowcolor{AnthDark}
  \color{white}\textbf{Algorithm} &
  \color{white}\textbf{Description} &
  \color{white}\textbf{Practical cost} \\
  \rowcolor{AnthLightGray!60}
  OPT       & Evict the page used furthest in the future & Impossible online \\
  LRU       & Evict least recently used page & High (full timestamp) \\
  \rowcolor{AnthLightGray!60}
  Clock     & Approximates LRU with a reference bit      & Low \\
  LFU       & Evict least frequently used page           & Medium \\
  \rowcolor{AnthLightGray!60}
  Linux (2Q)& Active/inactive list pair with promotion   & Low \\
\end{tabularx}
\end{center}

%% ─────────────────────────────────────────────────────────────────────────────
\section{Data Structures and Algorithms}
%% ─────────────────────────────────────────────────────────────────────────────

\subsection{Asymptotic Complexity}

\begin{Def}[Big-O notation]
  $f(n) = O(g(n))$ iff $\exists\, c > 0,\; n_0 > 0$ such that
  $f(n) \leq c \cdot g(n)$ for all $n \geq n_0$.
  It describes an \emph{asymptotic upper bound} — it says nothing about constant factors
  or small inputs.
\end{Def}

\begin{Note}
  $O$, $\Omega$, and $\Theta$ are distinct. When people say "the complexity \emph{is}
  $O(n \log n)$" they usually mean $\Theta(n \log n)$ — both upper and lower bounded.
  Sloppy usage propagates through a lot of textbooks; be precise in your own work.
\end{Note}

\subsection{Sorting algorithms — complexity summary}

\begin{center}
\begin{tabularx}{\linewidth}{L{2.8cm} C{2cm} C{2cm} C{2cm} C{1.6cm}}
  \rowcolor{AnthDark}
  \color{white}\textbf{Algorithm} &
  \color{white}\textbf{Best} &
  \color{white}\textbf{Average} &
  \color{white}\textbf{Worst} &
  \color{white}\textbf{Space} \\
  \rowcolor{AnthLightGray!60}
  Insertion sort & $O(n)$ & $O(n^2)$ & $O(n^2)$ & $O(1)$ \\
  Merge sort   & $O(n \log n)$ & $O(n \log n)$ & $O(n \log n)$ & $O(n)$ \\
  \rowcolor{AnthLightGray!60}
  Quick sort   & $O(n \log n)$ & $O(n \log n)$ & $O(n^2)$ & $O(\log n)$ \\
  Heap sort    & $O(n \log n)$ & $O(n \log n)$ & $O(n \log n)$ & $O(1)$ \\
  \rowcolor{AnthLightGray!60}
  Timsort      & $O(n)$ & $O(n \log n)$ & $O(n \log n)$ & $O(n)$ \\
\end{tabularx}
\end{center}

\subsection{Example: merge sort}

\begin{tcolorbox}[codebox, title=Merge sort — Python]
\begin{lstlisting}[language=Python]
def merge_sort(arr: list[int]) -> list[int]:
    """Recursive merge sort. Returns a new sorted list."""
    if len(arr) <= 1:
        return arr

    mid   = len(arr) // 2
    left  = merge_sort(arr[:mid])
    right = merge_sort(arr[mid:])
    return _merge(left, right)


def _merge(left: list[int], right: list[int]) -> list[int]:
    result, i, j = [], 0, 0
    while i < len(left) and j < len(right):
        if left[i] <= right[j]:
            result.append(left[i]); i += 1
        else:
            result.append(right[j]); j += 1
    return result + left[i:] + right[j:]
\end{lstlisting}
\end{tcolorbox}

\begin{Tip}[When to use Timsort]
  Python's built-in \ic{sorted()} and \ic{list.sort()} both use Timsort, which
  exploits existing \emph{runs} of sorted data in real inputs. For nearly-sorted
  sequences it degrades gracefully to $O(n)$. Prefer it over rolling your own
  sort for anything outside embedded or GPU contexts.
\end{Tip}

\subsection{Graph representations}

Two standard representations for a graph $G = (V, E)$:

\concept{Adjacency matrix}

An $|V| \times |V|$ matrix where entry $(i, j) = 1$ iff edge $(i, j) \in E$.
Space: $O(V^2)$. Edge lookup: $O(1)$. Preferred for dense graphs.

\concept{Adjacency list}

An array of $|V|$ lists, each containing the neighbours of vertex $i$.
Space: $O(V + E)$. Edge lookup: $O(\deg(v))$. Preferred for sparse graphs —
which is nearly every real-world network.

\begin{tcolorbox}[codebox, title=BFS — C]
\begin{lstlisting}[language=C]
#include <stdio.h>
#include <stdbool.h>
#include <stdlib.h>

#define MAX 128

void bfs(int adj[MAX][MAX], int n, int src) {
    bool visited[MAX] = {false};
    int  queue[MAX], front = 0, rear = 0;

    visited[src]  = true;
    queue[rear++] = src;

    while (front < rear) {
        int v = queue[front++];
        printf("%d ", v);
        for (int u = 0; u < n; u++) {
            if (adj[v][u] && !visited[u]) {
                visited[u]    = true;
                queue[rear++] = u;
            }
        }
    }
}
\end{lstlisting}
\end{tcolorbox}

%% ─────────────────────────────────────────────────────────────────────────────
\section{Networking}
%% ─────────────────────────────────────────────────────────────────────────────

\subsection{The TCP/IP Stack}

\begin{center}
\begin{tabularx}{\linewidth}{C{0.4cm} L{2.6cm} X L{3.8cm}}
  \rowcolor{AnthDark}
  \color{white}\textbf{\#} &
  \color{white}\textbf{Layer} &
  \color{white}\textbf{Responsibility} &
  \color{white}\textbf{Key protocols} \\
  \rowcolor{AnthLightGray!60}
  4 & Application  & User-facing services             & HTTP, DNS, SMTP, SSH \\
  3 & Transport    & End-to-end delivery, flow control & TCP, UDP, QUIC \\
  \rowcolor{AnthLightGray!60}
  2 & Internet     & Logical addressing, routing       & IP (v4/v6), ICMP \\
  1 & Link         & Framing, MAC addressing           & Ethernet, Wi-Fi \\
\end{tabularx}
\end{center}

\subsection{TCP three-way handshake}

\begin{enumerate}
  \item Client sends \ic{SYN} (seq = $x$).
  \item Server replies \ic{SYN-ACK} (seq = $y$, ack = $x+1$).
  \item Client sends \ic{ACK} (ack = $y+1$). Connection established.
\end{enumerate}

\begin{Warn}[SYN flood attacks]
  A SYN flood exploits the half-open connection state: an attacker sends many
  \ic{SYN} packets with spoofed source IPs, exhausting the server's connection
  backlog. Mitigated with \emph{SYN cookies}: the server encodes state into the
  \ic{ISN} instead of allocating a TCB, deferring allocation until the \ic{ACK}
  is verified.
\end{Warn}

\begin{Note}
  QUIC (RFC~9000) eliminates the TCP handshake latency for repeat connections via
  \emph{0-RTT resumption} and integrates TLS~1.3 natively — the entire handshake
  costs 1~RTT on first connection, 0~RTT on resumption.
\end{Note}

\end{document}